\documentclass[11pt,a4paper]{article}
\usepackage[margin=2.5cm]{geometry}
\usepackage{amsmath}
\usepackage{booktabs}
\usepackage{array}
\usepackage{hyperref}
\usepackage{parskip}
\usepackage{xcolor}
\usepackage{listings}
\usepackage{caption}
\usepackage{subcaption}

\lstset{
  basicstyle=\ttfamily\footnotesize,
  breaklines=true,
  frame=single,
  backgroundcolor=\color{gray!8},
}

\title{\textbf{Data Processing Pipeline} \\[0.5em]
       \large Gas Injury Detection --- MEMS and SPEC Sensors}
\author{Experiment Data}
\date{April 2026}

\begin{document}
\maketitle
\tableofcontents
\newpage

% ─────────────────────────────────────────────────────────────────────────────
\section{Overview}

Two independent processing pipelines clean the raw CSV recordings before
any modelling.  Each pipeline is implemented in a Jupyter notebook and
exports \texttt{.pkl} files to the \texttt{processed/} directory.

\begin{center}
\begin{tabular}{lll}
\toprule
Notebook & Sensors & Output files \\
\midrule
\texttt{mems\_filter.ipynb} & VOC, NH3, HCHO (metal-oxide) & \texttt{mems\_*.pkl} \\
\texttt{spec\_filter.ipynb} & H2S, EtOH (electrochemical) & \texttt{spec\_*.pkl} \\
\bottomrule
\end{tabular}
\end{center}

The datasets processed are:

\begin{center}
\begin{tabular}{llll}
\toprule
Key & Source file & Session & Notes \\
\midrule
SWEAT\_1a & \texttt{20260405-experiment/sweat.csv} & Apr 05 & Up to 4000\,s \\
SWEAT\_1b & \texttt{20260405-experiment/sweat.csv} & Apr 05 & 4000\,s onward; elapsed reset \\
BLOOD\_0  & \texttt{20260406-experiment/1.5\_blood.csv}   & Apr 06 & Full session \\
BLOOD\_1  & \texttt{20260406-experiment/1.5\_blood\_2.csv} & Apr 06 & Trimmed: $t \geq 1750$\,s \\
\bottomrule
\end{tabular}
\end{center}

\textbf{Sweat 0} (\texttt{20260404-experiment/sweat.csv}) is excluded from
training --- the NH3 pin was disconnected (miswired) for that session ---
but it is used as an out-of-sample validation target.

% ─────────────────────────────────────────────────────────────────────────────
\section{MEMS Pipeline (\texttt{mems\_filter.ipynb})}

\subsection{Data Loading}

Raw CSV is read with \texttt{wall\_time} parsed as a timestamp.
Elapsed time is computed as:
\[
  t_{\text{elapsed}} = \texttt{wall\_time} - \texttt{wall\_time}[0] \quad [\text{seconds}]
\]
An optional trim is applied: rows with $t < t_{\min}$ or $t \geq t_{\max}$ are
discarded before any processing.

For Sweat 1, the session is split at $t_{\text{reset}} = 4000$\,s (a recording
restart event):
\begin{itemize}
  \item \textbf{SWEAT\_1a}: $t < 4000$\,s --- contains sweat samples 2 and 3
        (introduced $\approx$3850\,s).
  \item \textbf{SWEAT\_1b}: $t \geq 4000$\,s --- elapsed reset to 0; contains
        sweat sample 1 (introduced $\approx$570\,s of sub-session time).
\end{itemize}

\subsection{Artifact Removal}

Known electrical or handling artifacts are patched \emph{before} smoothing by
replacing the affected time range with \texttt{NaN} and interpolating linearly
between the clean neighbours.  Per-dataset artifact windows:

\begin{center}
\begin{tabular}{llll}
\toprule
Dataset & Channel & $t_{\text{start}}$ (s) & $t_{\text{end}}$ (s) \\
\midrule
SWEAT\_1a & VOC, NH3, HCHO & 2950 & 3100 \\
SWEAT\_1b & HCHO           &  540 &  570 \\
BLOOD\_0  & NH3            & 1970 & 2030 \\
\bottomrule
\end{tabular}
\end{center}

BLOOD\_0 NH3 also has isolated zero readings interpolated before artifact
removal (sensor dropout artefacts).

\subsection{Smoothing --- Causal Exponential Moving Average}

After artifact patching, each MEMS channel is smoothed with a causal
Exponential Moving Average (EMA):
\[
  \hat{x}[n] = \alpha\, x[n] + (1 - \alpha)\, \hat{x}[n-1], \qquad
  \alpha = \frac{2}{\text{span} + 1}
\]

Default span is 30 samples ($\approx 30$\,s at 1\,Hz).  SWEAT\_1b NH3 uses
a larger span of 80 to suppress its higher noise level.

\textbf{Why causal EMA:} The model will run in real time on a microcontroller.
Using a non-causal filter (e.g.\ Savitzky--Golay, zero-phase Butterworth)
during training but a causal filter at inference would introduce a
train/inference mismatch.  Causal EMA is identical offline and online.

\subsection{Baseline and Sensor Drift Correction}

Metal-oxide sensors exhibit a slow resistive drift even in clean air.  A linear
trend is estimated from a \emph{pre-sample baseline window} and subtracted from
the full signal:

\begin{enumerate}
  \item Select rows in the baseline window $[t_0, t_1]$ (before sample
        introduction).
  \item Fit a degree-1 polynomial to the EMA-filtered signal in that window:
        \[
          \hat{y}(t) = a \cdot t + b
          \qquad \text{via least squares on } t \in [t_0, t_1]
        \]
  \item Subtract the extrapolated trend from the entire session:
        \[
          y_{\text{corr}}(t) = \max\!\bigl(y(t) - \hat{y}(t),\; 0\bigr)
        \]
        The $\max(\cdot, 0)$ clip prevents negative readings.
\end{enumerate}

Baseline windows per dataset:

\begin{center}
\begin{tabular}{lrr}
\toprule
Dataset & $t_0$ (s) & $t_1$ (s) \\
\midrule
SWEAT\_1a & 1500 & 2800 \\
SWEAT\_1b &    0 &  570 \\
BLOOD\_0  &    0 & 1900 \\
BLOOD\_1  & 1750 & 2440 \\
\bottomrule
\end{tabular}
\end{center}

\subsection{Export}

Drift-corrected channels \texttt{[elapsed\_s, voc, nh3, hcho, temp\_C, rh\_pct]}
are exported as \texttt{processed/mems\_\{key\}.pkl} for each dataset.

% ─────────────────────────────────────────────────────────────────────────────
\section{SPEC Pipeline (\texttt{spec\_filter.ipynb})}

\subsection{PPM Conversion}

Raw voltage readings from the SPEC electrochemical sensors are converted to PPM
using the transimpedance amplifier (TIA) gain and manufacturer sensitivity code:
\[
  m = S_{\text{code}} \times R_{\text{TIA}} \times 10^{-6} \quad [\text{V/ppm}]
\]
\[
  \text{ppm} = \max\!\left(\frac{V_{\text{gas}} + \delta V - V_{\text{ref}}}{m},\; 0\right) \times s
\]

where $\delta V$ is a per-sensor voltage offset and $s$ is a multiplicative
scale factor.  Parameters used:

\begin{center}
\begin{tabular}{lrrrrrr}
\toprule
Sensor & $S_{\text{code}}$ (nA/ppm) & $R_{\text{TIA}}$ (k$\Omega$) & $m$ (V/ppm) & $\delta V$ (V) & $s$ \\
\midrule
H2S  & 216.09 & 49.9  & $1.078 \times 10^{-2}$ & 0.000 & 1.5 \\
EtOH &  21.5  & 249.0 & $5.354 \times 10^{-3}$ & +0.100 & 1.0 \\
\bottomrule
\end{tabular}
\end{center}

\subsection{Outlier Detection and Handling}

Two independent detectors are run on the raw PPM signal; a sample flagged by
\emph{either} is treated as an outlier.

\subsubsection*{Rolling IQR}
\[
  \text{outlier if } x < Q_1 - k \cdot \text{IQR} \;\text{ or }\; x > Q_3 + k \cdot \text{IQR}
\]
Computed over a causal rolling window of 60 samples with $k = 2.0$.

\subsubsection*{Rolling Rate-of-Change (ROC)}
\[
  \text{outlier if } \frac{|\Delta x[n]| - \mu_{\text{ROC}}}{\sigma_{\text{ROC}}} > z_{\text{thresh}}
\]
where $\mu_{\text{ROC}}$ and $\sigma_{\text{ROC}}$ are the rolling mean and
standard deviation of $|\Delta x|$ over 60 samples, and $z_{\text{thresh}} = 3.0$.

Flagged samples are replaced by linear interpolation between the surrounding
clean values (forward/backward filled at boundaries).

\subsection{Smoothing --- Three-Stage Filter}

The outlier-handled signal is smoothed with a three-stage pipeline:

\begin{enumerate}
  \item \textbf{Savitzky--Golay (SG)} pre-smoother:
        window = 121 samples, polynomial order = 3.
        Removes rapid point noise while preserving peak shapes.

  \item \textbf{Causal Butterworth low-pass}:
        4th-order, $f_c = 0.005$\,Hz at $f_s = 1$\,Hz.
        Extracts the slow trend.  Implemented with \texttt{lfilter} and
        matched initial conditions (\texttt{lfilter\_zi}) to avoid edge
        transients.
\end{enumerate}

\noindent\textit{Note:} The SG stage is non-causal (used here for noise
pre-processing only); the Butterworth stage is causal.  At inference,
the SG step may be replaced by a running EMA of equivalent span.

\subsection{Boot-up Stabilization Detection}

After power-on, SPEC sensors drift substantially before settling.  Each sensor
uses a dedicated algorithm to detect the stabilization point.

\subsubsection*{H2S --- Dual Moving Average Convergence}

A short-term MA (60\,s) and long-term MA (300\,s) are computed on the
outlier-handled signal.  The sensor is considered stable when:
\[
  |\text{MA}_{\text{short}}(t) - \text{MA}_{\text{long}}(t)| < \theta
  \quad \text{for } \geq 300 \text{ consecutive samples}
\]
with $\theta = 0.05$\,ppm.  The first 1200 samples are always excluded (sensor
is guaranteed unstable during warmup).

\subsubsection*{EtOH --- Rolling Slope Threshold}

EtOH decays exponentially from a high boot value; the dual-MA gap closes too
slowly.  Instead, the three-stage filtered signal is used:
\[
  \text{stable when } \overline{|\Delta x|}_{w} < \theta_{\text{slope}}
  \quad \text{for } \geq 300 \text{ consecutive samples}
\]
where the rolling mean is computed over a long window and
$\theta_{\text{slope}}$ is a dataset-specific threshold.

\subsection{Stabilization Offset Correction}

The signal value at the detected stabilization point is used as a per-boot
baseline offset.  All readings are shifted so the baseline is zero:
\[
  y_{\text{corr}}(t) = \max\!\bigl(y(t) - y(t_{\text{stable}}),\; 0\bigr)
\]
If stabilization is not detected, the mean of the first 300 samples is used
as a fallback.

\subsection{Sweat 1 Session Split}

The Sweat 1 session (\texttt{20260405}) contains two distinct sub-sessions
separated by a recording restart at $t = 4000$\,s.  After offset correction,
the session is split:

\begin{itemize}
  \item \textbf{SPEC SWEAT\_1a}: $t \leq 4000$\,s
  \item \textbf{SPEC SWEAT\_1b}: $t > 4000$\,s, elapsed reset to 0
\end{itemize}

\subsection{Export}

Offset-corrected channels \texttt{[elapsed\_s, h2s\_ppm, etoh\_ppm]} are
exported as \texttt{processed/spec\_\{key\}.pkl} for each dataset.

% ─────────────────────────────────────────────────────────────────────────────
\section{Merge and Feature Engineering (\texttt{combined.ipynb})}

\subsection{Merge}

MEMS and SPEC pickle files are inner-joined on \texttt{elapsed\_s} to produce
one DataFrame per dataset with all seven channels:
\texttt{voc, nh3, hcho, h2s\_ppm, etoh\_ppm, temp\_C, rh\_pct}.

\subsection{Labelling}

Each row is assigned one of three class labels based on manually annotated
time windows (from experimental README):

\begin{center}
\begin{tabular}{rlll}
\toprule
Label & Class & Dataset & Window (s) \\
\midrule
0 & Baseline & All & Rows outside sample windows \\
1 & Sweat    & SWEAT\_1a & 2870 -- 3850 \\
1 & Sweat    & SWEAT\_1b &  570 -- 2330 \\
2 & Blood    & BLOOD\_0  & 1900 -- 3730 \\
2 & Blood    & BLOOD\_1  & 2440 -- 3800 \\
\bottomrule
\end{tabular}
\end{center}

\subsection{Helper Columns}

The following features are computed per row, causally (using only past data),
for each channel in \texttt{[voc, nh3, hcho, h2s\_ppm, etoh\_ppm, rh\_pct]}:

\begin{center}
\begin{tabular}{ll}
\toprule
Column & Description \\
\midrule
\texttt{\_roc}         & First difference $x[n] - x[n-1]$ \\
\texttt{\_acc}         & Second difference (diff of diff) \\
\texttt{\_roll\_mean\_\{w\}} & Causal rolling mean over $w \in \{15, 30, 60\}$ samples \\
\texttt{\_roll\_std\_\{w\}}  & Causal rolling std over $w$ samples \\
\texttt{\_roll\_roc\_\{w\}}  & Causal rolling mean of $|$ROC$|$ over $w$ samples \\
\bottomrule
\end{tabular}
\end{center}

\texttt{temp\_C} is excluded from feature columns; \texttt{rh\_pct} contributes
helper columns only (not its raw value).

\subsection{Windowed Feature Extraction}

Non-overlapping 60-second windows are slid separately over each
\emph{(session, label)} pair --- windows never cross label boundaries.
Within each window, four aggregates are computed per feature column:
\[
  \{\,\text{mean},\ \text{std},\ \text{max},\ \text{slope}\,\}
\]
where slope is the linear regression coefficient of the column against sample
index within the window.

\subsection{Class Balancing}

Baseline windows greatly outnumber sample windows.  The baseline set is
randomly undersampled to match the size of the larger sample class (sweat or
blood), producing a balanced \texttt{windows\_balanced} DataFrame.

% ─────────────────────────────────────────────────────────────────────────────
\section{Model --- Random Forest}

\subsection{Run 1}
\begin{itemize}
  \item 200 decision trees, \texttt{class\_weight="balanced"}
  \item All feature columns (raw aggregates + ROC + rolling mean/std/ROC + acceleration)
  \item Leave-One-Session-Out (LOSO) cross-validation via \texttt{GroupKFold(n\_splits=4)}
  \item Mean validation accuracy: 0.752, std: 0.253
  \item Fold 3 (held-out: blood\_0) collapsed to 0.333 --- blood\_0 EtOH max
        $\approx$14 ppm vs blood\_1 max $\approx$0.3 ppm; model had only seen blood\_1 for
        the blood class.
\end{itemize}

\subsection{Run 2}
\begin{itemize}
  \item 100 decision trees, \texttt{class\_weight="balanced"}
  \item Feature reduction: remove \texttt{\_roll\_mean\_*}, \texttt{\_roll\_std\_*},
        \texttt{\_acc} columns; keep raw aggregates and ROC columns only
  \item 80/20 stratified train/test split on \texttt{windows\_balanced}
  \item Per-class precision and recall reported on the held-out 20\%
\end{itemize}

Future sessions (new blood/sweat recordings) will serve as true out-of-sample
validation.

\end{document}
